\begin{center}\begin{tikzpicture}[font=\small,>=stealth,line join=round]
\draw[gray!40,rounded corners=3pt] (0,0) rectangle (14.2,8.5);\node at (7.1,7.9) {例二：奇次项可以组合，常数项必须为零};
\begin{scope}[shift={(4.2,3.8)},x=1.7cm,y=0.65cm]\draw[->,gray] (-1.8,0)--(1.8,0) node[right]{$x$};\draw[->,gray] (0,-4)--(0,4) node[above]{$y$};\draw[blue,thick,domain=-1.05:1.05,samples=150] plot(\x,{(\x)^5+(\x)^3+\x});\fill[orange!85!black] (1,3) circle(2.2pt) node[above right=5pt] {$(1,3)$};\fill[orange!85!black] (-1,-3) circle(2.2pt) node[below left=5pt] {$(-1,-3)$};\fill[orange!85!black] (0,0) circle(2.2pt) node[below right=5pt] {$O$};\end{scope}\node[align=left,anchor=west] at (8.1,4.2){图取 $f(x)=x^5+x^3+x$\\[11pt]每一项换 $-x$ 后都变号\\[11pt]任意实系数的奇次项之和\\[11pt]仍满足 $f(-x)=-f(x)$};
\node[text width=12.8cm,align=center] at (7.1,.65) {图取 $a=b=c=1,d=0$；一般情形看下方代数验证。};
\end{tikzpicture}\end{center}
